\section{Block architecture}
\label{sec:arch}

A block's network is Part I's energy model with the changes round 2 made; this section says what it is, what
has to change for blocks to be small and for lists and grids to be rules rather than tables, and what a block
must cost for a machine of them to run at frame rate.

\subsection{The energy block today}
\label{sec:energyblock}

A finite block scores a candidate successor $s'$ of $(s,a)$ by $E(s,a,s')=\Ej+\lambda_F\Ef$: the JEPA
prediction error between a predictor's output for $(s,a)$ and a target encoder's embedding of $s'$, plus the
factored per-field term whose shift kernel is equivariant along ordinal fields (Part I, \S4). It is trained
with InfoNCE over the candidate set, a hinge on spec-wrong negatives, and a variance floor, over every cell of
the spec; hole examples are oversampled twenty to one against written cells. The candidate set is the states
the invariants allow (\cref{sec:valid}), both in the loss and at the argmin, and their target embeddings are
computed once per set of weights, which is what makes a step on $3{,}125$ candidates take under a millisecond
on a CPU (\texttt{docs/paper2/timings.md}). A continuous block contrasts each true successor against the batch and sixteen near-miss perturbations,
and steps by sixty Adam iterations on the energy over its continuous fields from eleven seeds per discrete
assignment, exactly as Part I's continuous inference. The networks in \cref{tab:round2} have \minParamsK k to \maxParamsK k
parameters; three-layer GELU MLPs with $d=64$ throughout.

\subsection{Heads by field type}
\label{sec:heads}

The factored head's shift kernel made ``$+1$'' the same function at every value of an ordinal field, and Part
I's ablation showed that is what let a rule generalise from few cells. Round 2 trains on every cell, so
generalisation is not the goal, but the same property decides \emph{how many cells a rule needs} and therefore
how large a block may be. Reorder's residual errors under scarce data are mostly of one kind: a legal order in
which the moved item is not where it was dropped, because the head has no notion that the shift rule is the
same function for item $a$ as for item $e$.

We propose that the type of a field chooses the symmetry of its head, and that the spec's type system is where
that choice is made. An ordinal field keeps the shift kernel. A field that is one of several positions of
interchangeable items (reorder's $p_a,\dots,p_e$; kanban's slots) is a \emph{list} in disguise, and its head
should be permutation-equivariant: a per-item scorer that sees its own position, the moved item's old and new
positions, and a set summary of the others, with parameters shared across items~\citep{zaheer2017deepsets,lee2019settransformer}.
Under that head the rule ``if $\mathit{from}<p\le\mathit{to}$ then $p-1$'' is one function of three numbers;
the cells it needs are the distinct $(\mathit{from},\mathit{to},p)$ triples, of order $n^3$, not the
$n!\cdot n^2$ cells of the table. A field that is a position on a grid (fifteen, sokoban, rush hour) wants a
translation-equivariant head that sees displacements, not coordinates~\citep{cohen2016gcnn}. Enums keep the
plain head. The test is not generalisation for its own sake but size: a block whose rule is learned as a
function can be trained on a sample of its cells, which is the other exit from the enumeration wall.

\subsection{Candidates from a proposer}
\label{sec:futures}

The argmin over every valid state is exact and cheap at $10^3$ states and impossible at $10^{13}$. In the lab a
drag already restricts candidates to the futures a gesture can reach, derived by the oracle; a shipped block
cannot ask an oracle. But the factored head is a proposer: its per-field logits name the $k$ most likely values
of each field, and the product of those is a candidate set of $k^F$ states that the energy then ranks. For a
block whose rule changes one or two fields per action, $k=3$ covers the truth essentially always and the
candidate set is tens of states whatever the block's size. The plan is to make this the inference path for any
block over about $10^4$ states, keep the exhaustive argmin as the reference it is checked against on small
blocks, and report the proposer's recall separately so that a miss can never be mistaken for a wrong ranking.

\subsection{Compiled core, trained holes}
\label{sec:compiled}

Part I compiled every explicit transition of a finite spec into exact ReLU weights and trained only the holes
and the generalisation; its guarantee (the compiled term dominates any learned term where the spec speaks) is
the cleanest kind of exactness this programme has. Round 2 trained everything, because the Rust trainer does not
yet have the compiler. Per block the division is natural: compile the explicit transitions, train the holes
(hanoi's solver), and train the continuous parts. A compiled block has no seed to vary and nothing to score; the
digit of \cref{fig:odometer-spec} should be compiled, not trained, and the odometer experiment of
\cref{sec:demos} runs it both ways.

Blocks meet at ports, and a port needs a notion of compatibility. We take Part I's roadmap item seriously:
the energy is the type. Two blocks are compatible at a read port if the reader's expected domain and invariant
for the field match the owner's, checked statically at wiring time, and if their energies agree on the owner's
valid states for that field, checked by sampling. A block that is swapped for a retrained one must pass the
same check before the machine accepts it, which is the machine-level version of Part I's rule that the weights
that ship are the weights that were scored.

\subsection{Sizes and budgets}
\label{sec:budgets}

Blocks must stay small for three independent reasons: the trainer's cost in states times cells
(\cref{fig:cost}), the frame budget of a running machine, and the reviewer's ability to check a block
exhaustively. We fix targets rather than hopes. A finite block has at most about $2\times10^3$ valid states and
$2\times10^5$ parameters and steps in under $5$\,ms on a CPU core (under $1$\,ms measured on $3{,}125$
candidates); a continuous block steps in under $50$\,ms ($44$\,ms measured for the descent today;
\texttt{docs/paper2/timings.md}). At $60$ frames per second a tick has $16$\,ms; a cascade of ten finite
blocks fits in it on one core, and blocks that are not on the same cascade path step in parallel. Anything
that does not fit these budgets is two blocks.
